\section{Setup: soft gluon wave function and the evolution operator}\label{sec:setup}
In this section we summarise the framework of~\cite{Radhakrishnan:2026bxy}, on which the present
calculation is built, and collect all the coefficients of the evolution operator $\Omega$ that enter
the single inclusive cross section. The coefficients of order $g$ and $g^{2}$ were derived
in~\cite{Radhakrishnan:2026bxy}; the coefficient of order $g^{3}$ is obtained here. The light-cone
wave functions from which the coefficients are extracted are summarised in Appendix~\ref{app:lcwf}.

\subsection{Light-cone conventions}\label{subsec:conventions}
We use light-cone coordinates $x^{\pm}=(x^{0}\pm x^{3})/\sqrt{2}$ and denote transverse coordinates
by $\v x=(x^{1},x^{2})$. The projectile moves in the $+$ direction, $x^{+}$ plays the role of time,
and $p^{+}>0$ is the longitudinal momentum of a physical particle. The light-cone energy of an
on-shell gluon is $p^{-}=\v p^{2}/2p^{+}$, and we write
\begin{align}
&E_g(p)=\frac{\v p^2}{2p^+}\,,\qquad E_{gg}(p,q)=\frac{\v p^2}{2p^+}+\frac{\v q^2}{2q^+}\,,\notag\\
&d^3p\equiv dp^+d^2\v p\,,\qquad \delta^{(3)}(p-q)\equiv\delta(p^+-q^+)\,\delta^{(2)}(\v p-\v q)\,.
\end{align}
We work in the light-cone gauge $A^{+}=0$, in which the transverse components $A^{a}_{i}$,
$i=1,2$, are the dynamical fields. The gluon creation and annihilation operators satisfy
\begin{align}\label{eq:aadagger}
\big[a^a_i(k),a^{\dagger b}_j(p)\big]=(2\pi)^3\,\delta^{ab}\,\delta_{ij}\,\delta^{(3)}(k-p)\,.
\end{align}
For Fourier transforms in the transverse plane we use
\begin{align}\label{eq:FTconventions}
&a(k^+,\v k)=\int d^2z\,e^{-i\v k\cdot\v z}\,a(k^+,\v z)\,,\qquad
a^\dagger(k^+,\v k)=\int d^2z\,e^{i\v k\cdot\v z}\,a^\dagger(k^+,\v z)\,,\notag\\
&f(\v k)=\int d^2x\,e^{i\v k\cdot\v x}\,f(\v x)\,.
\end{align}
The opposite signs of the exponents for operators and for c-number functions are deliberate: with
this choice the colour charge density of Eq.~\eqref{eq:rho} below, $\rho^{a}(-\v p)=\int d^{2}x\,
e^{-i\v p\cdot\v x}\rho^{a}(\v x)$, is bilinear in $a^{\dagger}(\v k)\,a(\v k+\v p)$. In the transverse
coordinate representation the commutation relation becomes
\begin{align}\label{eq:commutation}
\big[a^a_i(k^+,\v z),a^{\dagger b}_j(p^+,\v z')\big]=2\pi\,\delta^{ab}\delta_{ij}\,\delta(k^+-p^+)\,\delta^{(2)}(\v z-\v z')\,.
\end{align}
One- and two-gluon states are normalised as
\begin{align}\label{eq:states}
&\ket{g^a_i(k^+,\v z)}=\frac{a^{\dagger a}_i(k^+,\v z)}{(2\pi)^{1/2}}\ket{0}\,,\notag\\
&\ket{g^a_i(k^+,\v z)\,g^b_j(p^+,\v z')}=\frac{a^{\dagger a}_i(k^+,\v z)\,a^{\dagger b}_j(p^+,\v z')}{2\pi}\ket{0}\,,
\end{align}
so that $\braket{g^b_j(p^+,\v z')|g^a_i(k^+,\v z)}=\delta^{ab}\delta_{ij}\,\delta(k^+-p^+)\,\delta^{(2)}(\v z-\v z')$
(see Appendix~\ref{app:states}).

\subsection{Light-cone QCD Hamiltonian}\label{subsec:LCQCD}
The light-cone gauge Hamiltonian is standard~\cite{Lublinsky:2016meo,Kovner:2005pe,Kovner:2005jc}.
We restrict to pure Yang--Mills theory, $N_f=0$: quark loops, which contribute to $\beta_0$ and to
$g\to q\bar q$ splitting at next-to-leading order, are not included, so that $\beta_0=11N_c/3$
throughout. The Hamiltonian splits into electric and magnetic parts,
\begin{align}\label{eq:LCQCD}
H_{\rm LC\,YM}=\int dx^{-}d^{2}\v x\left(\frac{1}{2}\,\Pi^{a}(x^{-},\v x)\,\Pi^{a}(x^{-},\v x)
+\frac{1}{4}\,F_{ij}^{a}(x^{-},\v x)\,F_{ij}^{a}(x^{-},\v x)\right),
\end{align}
with
\begin{equation}\begin{split}
&\Pi^{a}(x^{-},\v x)\equiv-\frac{1}{\partial^{+}}\big(D_{i}^{ab}\,\partial^{+}A_{i}^{b}\big),\\
&F_{ij}^{a}(x^{-},\v x)\equiv\partial_{i}A_{j}^{a}-\partial_{j}A_{i}^{a}-g\,f^{abc}A_{i}^{b}A_{j}^{c}.
\end{split}\end{equation}
Expanding in the gauge field gives $H_{\rm LC\,YM}=H_{0}+H_{\rm int}$. The free part,
\begin{align}\label{eq:free}
H_{0}=\int dx^{-}\,d^{2}\v x\,\frac{1}{2}\big(\partial_{i}A_{j}^{a}\big)^{2},
\end{align}
defines the free gluon Fock states built on the perturbative vacuum. The interaction part,
\begin{align}\label{eq:Hint}
H_{\rm int}=\int dx^{-}\,d^{2}\v x\,\bigg(&-g\,f^{abc}A_{i}^{b}A_{j}^{c}\,\partial_{i}A_{j}^{a}
+\frac{g^{2}}{4}\,f^{abc}f^{ade}A_{i}^{b}A_{j}^{c}A_{i}^{d}A_{j}^{e}\notag\\
&-g\,f^{abc}\big(\partial_{i}A_{i}^{a}\big)\frac{1}{\partial^{+}}\big(A_{j}^{b}\partial^{+}A_{j}^{c}\big)\notag\\
&+\frac{g^{2}}{2}\,f^{abc}f^{ade}\frac{1}{\partial^{+}}\big(A_{i}^{b}\partial^{+}A_{i}^{c}\big)\frac{1}{\partial^{+}}\big(A_{j}^{d}\partial^{+}A_{j}^{e}\big)\bigg),
\end{align}
generates transitions between gluonic sectors order by order in $g$. The first two terms are the
three- and four-gluon vertices, and the last two are the instantaneous interactions characteristic
of light-cone quantisation. We quantise canonically by expanding the transverse fields in creation
and annihilation operators,
\begin{align}\label{eq:gaugequant}
A_{i}^{a}(x)=\int_{0}^{\infty}\frac{dk^{+}}{2\pi}\int\frac{d^{2}\v k}{(2\pi)^{2}}\,\frac{1}{\sqrt{2k^{+}}}
\Big(a_{i}^{a}(k^{+},\v k)\,e^{-ik\cdot x}+a_{i}^{a\dagger}(k^{+},\v k)\,e^{ik\cdot x}\Big),
\end{align}
which brings the free Hamiltonian to the form
\begin{align}\label{eq:freecan}
H_{0}=\int_{0}^{\infty}\frac{dk^{+}}{2\pi}\int\frac{d^{2}\v k}{(2\pi)^{2}}\,\frac{\v k^{2}}{2k^{+}}\,
a_{i}^{a\dagger}(k^{+},\v k)\,a_{i}^{a}(k^{+},\v k).
\end{align}
The interaction Hamiltonian is written in terms of creation and annihilation operators in
Sec.~\ref{sec:Eikonalapprox}.

\subsection{Born--Oppenheimer formalism and the eikonal approximation}\label{sec:Eikonalapprox}
In the CGC, gluon modes are separated by their longitudinal momentum into valence modes (large
$k^{+}$) and soft modes~\cite{McLerran:1993ni,McLerran:1993ka}. The valence modes are the
energetic constituents of the hadron and act as colour sources, while the soft modes are treated as
a dynamical field. The two sectors are separated by longitudinal momentum cutoffs: soft gluons
occupy $\Lambda<k^{+}<\vee$ and valence modes $k^{+}\geq\vee$, where $\Lambda$ is the lower cutoff
of the soft sector and $\vee$ is the soft--valence boundary. The width of the soft window,
$\delta Y=\ln(\vee/\Lambda)$, is the rapidity interval over which the soft part of the wave function
is generated; modes with $k^{+}<\Lambda$ are not part of the projectile wave function and are
attributed to the target. Since the light-cone energy $p^{-}=\v p^{2}/2p^{+}$ of a valence mode is
much smaller than that of a soft mode, the valence modes evolve on much longer light-cone time
scales $x^{+}$; on the time scale of the soft dynamics they are frozen and act as static colour
sources. Accordingly, the gluon field is decomposed as
\begin{align}\label{eq:decompose}
A_{i}^{a}(x)=\underline{A}_{i}^{a}(x)+\overline{A}_{i}^{a}(x),
\end{align}
where $\underline{A}$ contains the soft modes, $\Lambda<k^{+}<\vee$, and $\overline{A}$ the
valence modes, $k^{+}\geq\vee$. Inserting this decomposition into $H_{\rm int}$ of
Eq.~\eqref{eq:Hint} produces soft--soft, soft--valence and valence--valence terms (explicit
structures in Appendix~\ref{sec:App2}). Since $k^{+}_{\rm soft}/k^{+}_{\rm valence}\ll1$, we work in
the eikonal approximation and drop terms suppressed by this ratio, such as
$\frac{1}{\partial^{+}}\big(\underline{A}_{j}^{b}\,\partial^{+}\overline{A}_{j}^{c}\big)$. The
interaction Hamiltonian then takes the form of Ref.~\cite{Lublinsky:2016meo},
\begin{align}\label{eq:expandH}
H_{\rm int}=H_{g}+H_{ggg}+H_{gggg}+H_{gg\text{-inst}}+H_{gggg\text{-inst}}+H_{V}.
\end{align}
The first five terms contain at least one soft field and generate the soft dynamics in the valence
background; $H_{V}$ contains only valence fields and does not affect the soft sector. In the
Born--Oppenheimer picture, the valence sector acts on the soft fields only through its colour charge
density,
\begin{align}\label{eq:rho}
\rho^{a}(-\v p)=-i f^{abc}\int_{\vee}^{\infty}\frac{dk^{+}}{2\pi}\int\frac{d^{2}\v k}{(2\pi)^{2}}\,
a_{j}^{\dagger b}(k^{+},\v k)\,a_{j}^{c}(k^{+},\v k+\v p),
\end{align}
which satisfies the colour algebra
\begin{align}\label{eq:rhoalgebra}
\big[\rho^{a}(\v k),\rho^{b}(\v p)\big]=i f^{abc}\rho^{c}(\v k+\v p)\,.
\end{align}
% CHECK: confirm the statement below (normal ordering, LM 2016) before submission.
With the Hamiltonian normal ordered, as in Ref.~\cite{Lublinsky:2016meo}, only $H_{g}$, $H_{ggg}$
and $H_{gg\text{-inst}}$ contribute to the coefficients of $\Omega$ needed here; $H_{gggg}$ and
$H_{gggg\text{-inst}}$ first contribute at higher order and are not written out. Substituting the mode
expansion, Eq.~\eqref{eq:gaugequant}, these three terms read (details in Appendix~\ref{sec:App2})
\begin{align}\label{eq:Hg}
H_{g}=\int_{\Lambda}^{\vee}\frac{dk^{+}}{2\pi}\int\frac{d^{2}\v k}{(2\pi)^{2}}\,
\frac{g\,\v k^{i}}{\sqrt{2}\,|k^{+}|^{3/2}}\Big[a_{i}^{a\dagger}(k^{+},\v k)\,\rho^{a}(-\v k)
+a_{i}^{a}(k^{+},\v k)\,\rho^{a}(\v k)\Big],
\end{align}
\begin{align}\label{eq:Hggg}
H_{ggg}=&-\int_{\Lambda}^{\vee}\frac{dk^{+}}{2\pi}\,\frac{dp^{+}}{2\pi}\,dq^{+}\int\frac{d^{2}\v k}{(2\pi)^{2}}\,
\frac{d^{2}\v p}{(2\pi)^{2}}\,d^{2}\v q\;\frac{i g f^{abc}}{2\sqrt{2k^{+}p^{+}q^{+}}}\notag\\
&\times\bigg[\Big(\v q^{i}-\frac{q^{+}}{p^{+}+q^{+}}\,\v k^{i}\Big)\,a_{i}^{a}(k)\,a_{j}^{b\dagger}(p)\,a_{j}^{c\dagger}(q)\,\delta^{(3)}(-k+p+q)\notag\\
&\quad+\Big(\v p^{i}+\v q^{i}+\frac{p^{+}+q^{+}}{q^{+}-p^{+}}\,\v k^{i}\Big)\,a_{i}^{a\dagger}(k)\,a_{j}^{b\dagger}(p)\,a_{j}^{c}(q)\,\delta^{(3)}(k+p-q)+{\rm h.c.}\bigg],
\end{align}
\begin{align}\label{eq:Hgginst}
H_{gg\text{-inst}}=&\int_{\Lambda}^{\vee}\frac{dp^{+}}{2\pi}\frac{dq^{+}}{2\pi}\int\frac{d^{2}\v p}{(2\pi)^{2}}\frac{d^{2}\v q}{(2\pi)^{2}}\,
\frac{i g^{2} f^{abc}}{\sqrt{p^{+}q^{+}}}\notag\\
&\times\bigg(\frac{q^{+}\,\rho^{a}(-\v p-\v q)}{2(p^{+}+q^{+})^{2}}\,a_{i}^{b\dagger}(p^{+},\v p)\,a_{i}^{c\dagger}(q^{+},\v q)-{\rm h.c.}\notag\\
&\qquad-\frac{(p^{+}+q^{+})\,\rho^{a}(-\v p+\v q)}{2(p^{+}-q^{+})^{2}}\,a_{i}^{b\dagger}(p^{+},\v p)\,a_{i}^{c}(q^{+},\v q)\bigg).
\end{align}
